% Standalone source-aware packet salvaged from the rejected P01 converse.
% Do not infer or include the rejected global converse from this proposition.
% No novelty or priority claim is intended.

\paragraph{Source scope.}
The proof uses the following statements of Chambert--Loir--Ducros, in the
25 July 2025 version of \emph{Formes diff\'erentielles r\'eelles et courants
sur les espaces de Berkovich}.
\begin{enumerate}
\item Proposition~8.7.2 says that, if $T_k\to T$ are positive closed
currents, if $u_{0,k}\to u_0$ locally uniformly are continuous, and if
$u_{j,k}\to u_j$ locally uniformly are continuous psh functions for
$1\leq j\leq q$, then
\[
 u_{0,k}\,\mathrm d'\mathrm d''u_{1,k}\wedge\cdots\wedge
 \mathrm d'\mathrm d''u_{q,k}\wedge T_k
 \longrightarrow
 u_0\,\mathrm d'\mathrm d''u_1\wedge\cdots\wedge
 \mathrm d'\mathrm d''u_q\wedge T.
\]
The outer factor need only be continuous; it need not be psh.
\item In the construction in the proof of Proposition~16.9.1, a smooth
convex regularized maximum $M_\varepsilon$ is chosen nondecreasing in every
variable, with
\[
 \max_i s_i\leq M_\varepsilon(s)\leq\max_i s_i+\varepsilon,
\]
and with the deletion property: an entry at least $2\varepsilon$ below the
maximum does not affect $M_\varepsilon$.  Consequently the regularized
maximum of logarithmic norms is smooth and psh even across the zero loci of
the analytic functions which occur as entries.
\item Corollary~18.6.2 says that, on a boundaryless topologically separated
equidimensional analytic space, for analytic $f_1,\ldots,f_d$ and $r>0$,
\[
 \left(\mathrm d'\mathrm d''
  \log\max\{r,|f_1|,\ldots,|f_d|\}\right)^d
 =\delta_{f^{-1}(\eta_r)},
\]
where $\eta_r$ is the Gauss point of the polydisc of polyradius
$(r,\ldots,r)$ and the fibre current has its analytic multiplicities.
\item Proposition~18.7.3 says, with the relative dimension $n-d$ used in
its proof, that if $f=(f_1,\ldots,f_d):X\to\mathbf A^d$ is purely of
relative dimension $n-d$, the $f_i$ form a regular sequence, and
$Y=f^{-1}(0)$, then every compactly supported $G$-form $\omega$ of bidegree
$(n-d,n-d)$ with locally bounded tropically Borel coefficients satisfies
\[
 \langle\delta_{f^{-1}(t)},\omega\rangle
 =\langle\delta_Y,\omega\rangle
\]
for $t$ sufficiently close to the origin.  The displayed statement in the
source says ``relative dimension $d$'', but the next sentence, the
dimension argument, and the proof all use $n-d$; this packet uses that
proof-correct value.
\end{enumerate}

\begin{proposition}
Let $K$ be a complete nontrivially valued non-Archimedean field.  Let $X$ be
a good separated boundaryless strictly $K$-analytic space, pure of dimension
$n\geq2$, and let $h:X\to\mathbf R$ be a continuous $\mathbf R$-PL
function such that
\[
                         \mathrm d'\mathrm d''[h]\geq0.
\]
Let $V\subset X$ be an open analytic subspace and let
$f_1,\ldots,f_{n-1}\in\mathcal O_X(V)$ be a regular sequence.  Suppose that
$f=(f_1,\ldots,f_{n-1}):V\to\mathbf A^{n-1}$ is purely of relative
dimension one, and put $C=f^{-1}(0)$, endowed with its analytic-cycle
multiplicities.

For every nonnegative smooth compactly supported function $\psi$ on $V$,
\[
 \bigl\langle\mathrm d'\mathrm d''[h|_C],\psi|_C\bigr\rangle\geq0.
\]
In particular, suppose that $z\in C$, that
$\mathrm d'\mathrm d''[h|_C]$ is a discrete signed measure on an open
neighbourhood $W$ of $z$ in $C$, and that there is a nonnegative smooth
compactly supported function $\psi$ on $V$ which equals one near $z$ and
satisfies
\[
 \operatorname{Supp}(\psi)\cap C\subset W,
 \qquad
 \operatorname{Supp}(\psi)\cap
 \operatorname{Supp}\bigl(\mathrm d'\mathrm d''[h|_C]\bigr)=\{z\}.
\]
Then the coefficient of $\mathrm d'\mathrm d''[h|_C]$ at $z$ is
nonnegative.
\end{proposition}

\begin{proof}
Set $d=n-1$.  The open analytic subspace $V$ is again boundaryless,
separated, and pure of dimension $n$.  For $r>0$, put
\[
 u_r=\log\max\{r,|f_1|,\ldots,|f_d|\}.
\]
Using the regularized maximum from the proof of Chambert--Loir--Ducros,
Proposition~16.9.1, set
\[
 u_{r,\varepsilon}
 =M_\varepsilon(\log r,\log|f_1|,\ldots,\log|f_d|).
\]
At a point where some $f_i$ vanishes, $\log|f_i|$ tends to $-\infty$ and is
at least $2\varepsilon$ below the finite entry $\log r$ on a neighbourhood;
the deletion property makes $M_\varepsilon$ independent of that entry
there.  Thus $u_{r,\varepsilon}$ is globally well defined and smooth on
$V$.  Convexity and coordinatewise monotonicity make it psh.  Moreover
\[
 u_r\leq u_{r,\varepsilon}\leq u_r+\varepsilon,
\]
so $u_{r,\varepsilon}\to u_r$ uniformly on $V$.

Write
\[
 A_{r,\varepsilon}
   =(\mathrm d'\mathrm d''u_{r,\varepsilon})^d.
\]
This is a smooth positive closed $(d,d)$-form.  Hence
$\psi A_{r,\varepsilon}$ is a smooth positive compactly supported
$(n-1,n-1)$-form, the correct test bidegree for the $(1,1)$-current
$\mathrm d'\mathrm d''[h]$.  Positivity gives
\begin{equation}
\label{eq:p01-lci-positive}
 0\leq
 \left\langle\mathrm d'\mathrm d''[h],
              \psi A_{r,\varepsilon}\right\rangle.
\end{equation}
The definition of current differentials has no sign in this twice-applied
degree-zero case.  Since $A_{r,\varepsilon}$ is both
$\mathrm d'$-closed and $\mathrm d''$-closed, the Leibniz rule gives
\begin{align}
 \left\langle\mathrm d'\mathrm d''[h],
              \psi A_{r,\varepsilon}\right\rangle
 &=\left\langle[h],
       \mathrm d'\mathrm d''(\psi A_{r,\varepsilon})\right\rangle\notag\\
 &=\int_V h\,\mathrm d'\mathrm d''\psi\wedge A_{r,\varepsilon}.
\label{eq:p01-lci-parts}
\end{align}

Apply Chambert--Loir--Ducros, Proposition~8.7.2 with the constant sequence
of positive closed currents $T_\varepsilon=\delta_V$, with all $d$ psh
factors equal to $u_{r,\varepsilon}$, and with the constant continuous outer
factor $h$.  It yields weak convergence
\[
 hA_{r,\varepsilon}
 \longrightarrow
 h(\mathrm d'\mathrm d''u_r)^d.
\]
Pairing this convergence with the fixed smooth compactly supported
$(1,1)$-form $\mathrm d'\mathrm d''\psi$ and using
\eqref{eq:p01-lci-parts} gives
\begin{equation}
\label{eq:p01-lci-bt-limit}
 \lim_{\varepsilon\to0}
 \left\langle\mathrm d'\mathrm d''[h],
              \psi A_{r,\varepsilon}\right\rangle
 =\left\langle
   h(\mathrm d'\mathrm d''u_r)^d,
   \mathrm d'\mathrm d''\psi\right\rangle.
\end{equation}
By Corollary~18.6.2, the right-hand side is
\[
 \left\langle\delta_{f^{-1}(\eta_r)},
                   h\,\mathrm d'\mathrm d''\psi\right\rangle.
\]

We check the coefficient hypothesis needed for Proposition~18.7.3 rather
than treating it as automatic.  On a $G$-cover, an $\mathbf R$-PL function
has the form
\[
                   h=\sum_{j=1}^N a_j\log|g_j|
\]
with $a_j\in\mathbf R$ and analytic units $g_j$.  It is therefore the
pullback, under the corresponding moment map, of a continuous affine
function.  Hence $h$ is $G$-locally tropically continuous.  Multiplication
of the smooth $G$-form $\mathrm d'\mathrm d''\psi$ by $h$ consequently
produces a $G$-form with tropically continuous, in particular locally
bounded tropically Borel, coefficients.  Its support is contained in the
compact set $\operatorname{Supp}(\psi)$, and its bidegree is $(1,1)$,
which equals $(n-d,n-d)$.

The assumed pure relative dimension is $1=n-d$, so Proposition~18.7.3
applies.  As $r\to0$, the Gauss point $\eta_r$ tends to the origin of
$\mathbf A^d$.  Thus for all sufficiently small $r>0$,
\begin{align}
 \left\langle\delta_{f^{-1}(\eta_r)},
                   h\,\mathrm d'\mathrm d''\psi\right\rangle
 &=\left\langle\delta_C,
                   h\,\mathrm d'\mathrm d''\psi\right\rangle\notag\\
 &=\left\langle[h|_C],
                   \mathrm d'\mathrm d''(\psi|_C)\right\rangle\notag\\
 &=\left\langle\mathrm d'\mathrm d''[h|_C],
                   \psi|_C\right\rangle.
\label{eq:p01-lci-specialize}
\end{align}
Here $\delta_C$ and the restriction current use the same analytic-cycle
multiplicities.  Combining \eqref{eq:p01-lci-positive},
\eqref{eq:p01-lci-bt-limit}, and \eqref{eq:p01-lci-specialize} proves the
first assertion.

Under the hypotheses of the last assertion, the support of $\psi|_C$ lies
inside the neighbourhood where the current is a discrete signed measure,
and it meets the support of that measure only at $z$.  Since
$\psi|_C=1$ near $z$, the pairing just proved nonnegative is exactly the
coefficient of the atom at $z$.  This proves the last assertion.
\end{proof}

\paragraph{Scope firewall.}
This proposition does not assert that every nonnegative smooth function on
$C$ extends to a nonnegative smooth function on $V$, and therefore it does
not by itself prove positivity of the full restricted current against all
curve test functions.  It also contains no degree computation and implies
none of the rejected global converse.
